% from ega4-2-fr.tex
\oldpage[IV-2]{73}

En outre, si $X$ est localement de type fini sur $k$, ces conditions sont aussi
équivalentes à la suivante~:
\begin{enumerate}
\item[d)] Pour toute extension $k'$ de $k$ (ou pour une extension
algébriquement close $k'$ de $k$), $\mathscr F'$ est un
$\mathcal O_{X'}$-Module réduit.
\end{enumerate}
\end{proposition}

\begin{proof}
Il est clair que a) entraîne b). La condition b) entraîne d'abord que
$\mathscr F$ est réduit~: cela signifie (3.2.3) que si $Y$ est le
sous-préschéma fermé de $X$ défini par l'Idéal cohérent $\mathscr J$ de
$\mathcal O_X$, annulateur de $\mathscr F$, $Y$ est réduit, et il y a un
$\mathcal O_Y$-Module cohérent sans torsion $\mathscr G$ de rang~1 sur toute
composante irréductible de $Y$, tel que $j_*(\mathscr G)=\mathscr F$,
$j:Y\to X$ étant l'injection canonique. Or, pour toute extension $k'$ de $k$,
l'annulateur de $\mathscr F'$ est $\mathscr J'=\mathscr J\otimes_k k'$, le
morphisme projection $X'\to X$ étant plat (2.1.11)~; le sous-préschéma fermé de
$X'$ défini par $\mathscr J'$ est $Y'=Y\otimes_k k'$, et si
$j':Y'\to X'$ est l'injection canonique, et
$\mathscr G'=\mathscr G\otimes_k k'$, on a $\mathscr F'=j'_*(\mathscr G')$.
Notons aussi que, pour toute extension $k'$ de $k$ telle que $X'$ soit
localement noethérien, $\mathscr F'$ est sans cycle premier associé immergé
(4.2.7). De plus, tout point maximal $y'$ de $Y'$ est au-dessus d'un point
maximal $y$ de $Y$~; comme $\mathscr G_y$ est isomorphe à
$\mathcal O_{Y,y}$, $\mathscr G'_{y'}$ est isomorphe à
$\mathcal O_{Y',y'}$. On voit donc (par (3.2.3)) qu'il revient au même de dire
que $\mathscr F'$ est réduit ou que le préschéma $Y'$ est réduit, lorsque $X'$
est localement noethérien. Le fait que b) entraîne c) et que c) entraîne
a) est donc conséquence de (4.6.1, d) et b))~; dans a), l'hypothèse que
$k'$ est une extension de type fini de $k$ n'est faite que pour assurer que
$X'$ est localement noethérien. Lorsque $X$ est localement de type fini, $X'$
est localement noethérien pour toute extension $k'$ de $k$, ce qui prouve dans
ce cas l'équivalence de d) et des autres conditions.
\end{proof}

\begin{definition}[4.6.17]
\label{IV.2.4.6.17-fr}
Lorsque $\mathscr F$ vérifie les conditions équivalentes de (4.6.16), on dit
que $\mathscr F$ est \emph{géométriquement réduit} sur $k$, ou \emph{séparable}
sur $k$.
\end{definition}

\begin{proposition}[4.6.18]
\label{IV.2.4.6.18-fr}
Soient $k$ un corps~; $X$ un $k$-préschéma localement noethérien, $\mathscr F$
un $\mathcal O_X$-Module cohérent. Les conditions suivantes sont
équivalentes~:
\begin{enumerate}
\item[a)] Pour toute extension de type fini $k'$ de $k$, si l'on pose
$X'=X\otimes_k k'$, alors $\mathscr F'=\mathscr F\otimes_k k'$ est un
$\mathcal O_{X'}$-Module intègre (3.2.4).
\item[b)] Pour toute extension finie $k'$ de $k$, $\mathscr F'$ est un
$\mathcal O_{X'}$-Module intègre.
\item[c)] $\mathscr F$ est réduit (ou intègre), et si $\mathscr J$ est
l'Idéal de $\mathcal O_X$ annulateur de $\mathscr F$, le sous-préschéma fermé
défini par $\mathscr J$ est géométriquement intègre.
\end{enumerate}
En outre, si $X$ est localement de type fini sur $k$, ces conditions sont aussi
équivalentes à la suivante~:
\begin{enumerate}
\item[d)] Pour toute extension $k'$ de $k$ (ou pour une extension
algébriquement close $k'$ de $k$), $\mathscr F'$ est un
$\mathcal O_{X'}$-Module intègre.
\end{enumerate}
\end{proposition}

\begin{proof}
Les conditions a), b) et c) entraînent que $\mathscr F$ est intègre, donc
$\operatorname{Ass}(\mathscr F)$ réduit à un seul point $x$, qui est le point
générique de $\operatorname{Supp}(\mathscr F)$. Pour toute extension $k'$ de
$k$ pour laquelle $X'$ est localement noethérien, les cycles premiers associés
à $\mathscr F'$ sont donc non immergés, et sont les points maximaux de
$\operatorname{Supp}(\mathscr F')$, qui sont tous au-dessus de $x$. D'autre
part, en vertu de (4.6.16), chacune des conditions a), b), c) entraîne que
$\mathscr F$ est géométriquement réduit~; il est trivial que a) entraîne b),
et b) entraîne que $\operatorname{Supp}(\mathscr F)$ est géométriquement
irréductible (4.5.9), donc c), puisque le sous-préschéma $Y$ défini

